\section{Implementation}
\label{sec:implementation}
We reimplemented Eclat from Zaki's original paper~\cite{zaki2000} in Julia, using only
the standard library and no frequent itemset mining (FIM) library. The package has two
variants that return the same result: a baseline that follows the paper closely, and an
optimized variant. This section explains the algorithm and why it is correct, how we test
it, which optimizations we added, and how much each one improved on the baseline.

\subsection{Vertical layout and tidsets}
A database holds $N$ transactions. Each transaction gets a transaction identifier (tid),
the index $1, 2, \dots, N$ in the order it is read. Apriori stores the database in a
horizontal layout, where each transaction carries its list of items. Eclat uses a
vertical layout, where each item $x$ carries its tidset $t(x)$, the set of tids of the
transactions that contain $x$. For an itemset $X$,
\begin{equation}
  t(X) = \{\, i \in \{1,\dots,N\} : X \subseteq T_i \,\}, \qquad
  \sigma(X) = |t(X)|,
  \label{eq:support}
\end{equation}
where $T_i$ is the $i$-th transaction and $\sigma(X)$ is the support of $X$, the number
of transactions that contain $X$. An itemset $X$ is frequent when
$\sigma(X) \ge \mathit{minsup}$, with $\mathit{minsup}$ the minimum support threshold
counted in transactions. An empty transaction still has a tid and still counts toward
$N$, but it appears in no tidset.

The baseline stores each tidset as an increasing \texttt{Vector\{Int\}}, called a
tid-list. The function \texttt{to\_class} reads the database once and appends tids in
increasing order, so every tid-list is already sorted and never needs sorting again.

The main advantage of the vertical layout is the intersection. A transaction contains
$P \cup \{a, b\}$ if and only if it contains both $P \cup \{a\}$ and $P \cup \{b\}$, so
\begin{equation}
  t(P \cup \{a, b\}) = t(P \cup \{a\}) \cap t(P \cup \{b\}),
  \label{eq:intersect}
\end{equation}
where $P$ is an itemset and $a, b$ are two items not in $P$. The support of the new
itemset therefore comes from two tidsets that already exist, without another scan of the
database.

\subsection{Prefix-based equivalence classes and Bottom-Up search}
Zaki partitions the search space into prefix-based equivalence
classes~\cite{zaki2000}. The class with prefix $P$ holds the members
$(b, t(P \cup \{b\}))$, one for each item $b$ greater than every item of $P$. Two itemsets
in the same class differ only in their last item, so intersection~\eqref{eq:intersect}
is valid only between members of the same class.

The function \texttt{bottom\_up!} implements the paper's Bottom-Up procedure. For each
member $a$ of the class with prefix $P$, it records $P \cup \{a\}$ and its support,
intersects $t(P \cup \{a\})$ with the tidset of every member $b$ after $a$, keeps the
frequent results as the child class with prefix $P \cup \{a\}$, and searches that child
class before moving to the next member. The root class has the empty prefix and holds the
frequent items. The paper allows a depth-first search (DFS) or a breadth-first search over
the lattice of classes and reports results breadth-first. We chose depth-first so that
memory holds only one class per level of the current path. The paper also counts the
frequent 2-itemsets in advance with a triangular array over the horizontal database. The
baseline starts intersecting from single items so that one procedure handles every
level, and the optimized variant brings this counting step back in
Section~\ref{sec:optimized}. Table~\ref{tab:mapping} maps the parts of the paper to the
source code.

\begin{table}[t]
\centering
\small
\caption{Mapping from the paper~\cite{zaki2000} to the source code in \texttt{src/algorithm/}}
\label{tab:mapping}
\begin{tabular}{>{\raggedright\arraybackslash}p{0.37\columnwidth}>{\raggedright\arraybackslash}p{0.53\columnwidth}}
\hline
Paper & Baseline; optimized variant \\
\hline
Item tid-lists (Section 4.1) & \texttt{to\_class}; \texttt{bitset\_tidsets} \\
Prefix classes (Section 4.2) & vector of (item, tid-list) pairs; vector of
\texttt{ClassMember\{T\}} \\
Bottom-Up (Sections 4.3.1 and 5.2, Fig.~12) & \texttt{bottom\_up!}; \texttt{eclat\_search!} \\
Intersecting two tid-lists & \texttt{intersect\_tidlists}; \texttt{intersect\_bitsets!} and
\texttt{intersect\_lists!} \\
Counting 2-itemsets with a triangular array (Section 5.1) & not used; \texttt{count\_pairs} \\
\hline
\end{tabular}
\end{table}

Table~\ref{tab:trace} traces the algorithm on the database \texttt{baskets.txt}, which has
four transactions $T_1 = \{1,2,3\}$, $T_2 = \{1,2\}$, $T_3 = \{1,3\}$, $T_4 = \{2,3\}$,
with $\mathit{minsup} = 2$. The initial tidsets are $t(1) = \{1,2,3\}$,
$t(2) = \{1,2,4\}$, and $t(3) = \{1,3,4\}$. The algorithm records exactly six frequent
itemsets, and it prunes $\{1,2,3\}$ because only $T_1$ contains it.

\begin{table}[t]
\centering
\caption{Bottom-Up trace on \texttt{baskets.txt} with $\mathit{minsup} = 2$}
\label{tab:trace}
\begin{tabular}{llll}
\hline
Prefix & Intersection & Tidset & Result \\
\hline
$\emptyset$ & -- & $t(1) = \{1,2,3\}$ & record $\{1\}$, $\sigma = 3$ \\
$\{1\}$ & $t(1) \cap t(2)$ & $\{1,2\}$ & keep $2$ \\
$\{1\}$ & $t(1) \cap t(3)$ & $\{1,3\}$ & keep $3$ \\
$\{1\}$ & -- & $t(12) = \{1,2\}$ & record $\{1,2\}$, $\sigma = 2$ \\
$\{1,2\}$ & $t(12) \cap t(13)$ & $\{1\}$ & prune, $1 < 2$ \\
$\{1\}$ & -- & $t(13) = \{1,3\}$ & record $\{1,3\}$, $\sigma = 2$ \\
$\emptyset$ & -- & $t(2) = \{1,2,4\}$ & record $\{2\}$, $\sigma = 3$ \\
$\{2\}$ & $t(2) \cap t(3)$ & $\{1,4\}$ & keep $3$ \\
$\{2\}$ & -- & $t(23) = \{1,4\}$ & record $\{2,3\}$, $\sigma = 2$ \\
$\emptyset$ & -- & $t(3) = \{1,3,4\}$ & record $\{3\}$, $\sigma = 3$ \\
\hline
\end{tabular}
\end{table}

\subsection{Correctness}
\label{sec:correctness}
Every class that \texttt{bottom\_up!} receives satisfies four invariants: (i) each
tid-list is increasing and has no duplicates; (ii) the members are sorted by item and
every member item is greater than every item of the prefix; (iii) every member is
frequent; (iv) the tid-list of member $b$ equals $t(P \cup \{b\})$. The root class
satisfies all four when it is built. A child class takes the members after $a$, so it
keeps (i) and (ii). The size check before a member is added gives (iii), and
identity~\eqref{eq:intersect} gives (iv).

Three properties follow. First, every recorded support is exact, by (iv)
and~\eqref{eq:support}. Second, no itemset is recorded twice: by (ii) every key
$(x_1, \dots, x_k)$ is strictly increasing, and it is recorded only when the class with
prefix $(x_1, \dots, x_{k-1})$ reaches $x_k$, and each class is built only once. Third,
the result is complete because support is anti-monotone: if $X \subseteq Y$ then
$\sigma(X) \ge \sigma(Y)$, so pruning an infrequent itemset never loses a frequent
superset. Induction on length shows that the class with prefix $(x_1, \dots, x_{j-1})$
contains every item $y > x_{j-1}$ for which $(x_1, \dots, x_{j-1}, y)$ is frequent, so
every frequent itemset is recorded. The repository file \texttt{docs/eclat.md} gives the
full proof.

\subsection{Complexity}
Let $L$ be the total number of items over all transactions, $m$ the number of distinct
items, and $C$ the number of intersections. Building the vertical layout takes $O(L)$
expected time for the dictionary, $O(m \log m)$ for sorting, and $O(L)$ memory. One
intersection costs $O(|t(P \cup \{a\})| + |t(P \cup \{b\})|)$, at most $O(N)$. A class
with $k$ members performs $k(k-1)/2$ intersections, so the total time is
\begin{equation}
  O(L + m \log m + C \cdot N),
  \label{eq:cost}
\end{equation}
where $C$ is the total number of intersections over all classes visited. Because the
search is depth-first, memory holds only the classes on the current path, one class per
level, each with at most $m$ members of at most $N$ tids. The output can grow
exponentially: a frequent $k$-itemset has $2^k - 1$ non-empty subsets, and all of them
are frequent.

\subsection{Interface, SPMF format, and command line}
The two applications and the experiments call only four public functions.
\texttt{read\_transactions} reads an SPMF file (one transaction per line, items separated
by spaces). \texttt{support\_threshold} converts $\mathit{minsup}$ to a transaction
count. \texttt{mine} returns a \texttt{SupportMap} dictionary from itemset to support.
\texttt{write\_itemsets} writes one itemset per line in the form
\texttt{1 2 \#SUP: 2}. A real $\mathit{minsup}$ in $(0, 1]$ is a fraction and is
rounded up to $\lceil \mathit{minsup} \cdot N \rceil$ transactions; an integer
$\mathit{minsup}$ is an absolute transaction count. The \texttt{variant} argument
selects \texttt{:baseline} or \texttt{:optimized}, so each variant can be run on its own
and compared on the same input.

The command-line program \texttt{demo.jl} takes an input path, \texttt{--output},
exactly one of the two thresholds \texttt{--minsup} (a fraction) or
\texttt{--min-count} (a transaction count), and \texttt{--variant}, for example
\texttt{julia --project=. demo.jl --minsup 0.5 --output out.txt data/toy/baskets.txt}.
Invalid input stops the program with exit code 1 and a message that names the offending
line. The writer validates the whole result before it opens the file, so a failed run
leaves an existing output file untouched. Itemsets are written in lexicographic order of
their item codes, so the output files of the two variants are identical byte for byte.

\subsection{Automated tests}
\label{sec:tests}
The command \texttt{julia --project=. test/runtests.jl} runs 1\,091 checks, and all of
them pass on Julia 1.13.2. Each mining test compares the whole \texttt{SupportMap} and
also the number of itemsets recorded, because the dictionary would hide an itemset
recorded twice. Both variants run the following groups of tests.

\begin{itemize}
\item Hand-derived results: eight small databases, including an empty database, a
database with no frequent itemset, duplicate transactions, thresholds just above and
just below an itemset's support, and an image database with empty transactions.
\item Comparison with SPMF's Eclat~\cite{spmf2016}: eight runs on seven small databases
and one run on T25I10D10K (9\,976 transactions, 929 items), where SPMF reports 14\,782
itemsets of 1 to 13 items at $\mathit{minsup} = 0.008$. Both variants match 100\% on all
nine runs, with no missing, extra, or wrong-support itemset.
\item Comparison with a brute-force counter that enumerates every subset: 10 random
databases generated with seed 42 for each variant. The optimized variant also runs eight
databases with 0, 1, 63, 64, 65, 127, 128, and 129 transactions, which sit around the
boundary of one 64-bit bitset word. Half of them contain groups of items that always
occur together, to produce perfect extensions. Each database runs under 18
configurations that switch each technique of Section~\ref{sec:optimized} on and off.
\item Type stability: \texttt{Test.@inferred} on \texttt{mine}, \texttt{mine\_packed},
and the search functions of both variants.
\item Reading, writing, and the command line: the byte-level reader is checked against a
line-based parser on edge-case inputs (\texttt{Int} overflow, unexpected characters,
blank lines, Windows line endings) and on every small database. The writer is compared
byte for byte with a reference formatter, and argument errors and process exit codes are
checked.
\end{itemize}

The brute-force counter is only a test tool. Outside the test suite, every measurement
campaign in Section~\ref{sec:measure} compares each output file with an independent run
of SPMF's Eclat after timing ends. The final campaign produced 820 output files on the
five benchmark databases (615 from the optimized variant and 205 from PrePost+), and all
820 match 100\%.

\subsection{Optimized variant}
\label{sec:optimized}
The optimized variant keeps the Bottom-Up procedure and changes three things: how
tidsets are stored, which intersections are skipped, and how the result is written. Each
change keeps the output exact, and the proofs are in \texttt{docs/eclat-optimized.md}.
Start-up, reading, and writing live in shared code, so both variants benefit from the
changes there.

\subsubsection{Bitset tidsets and search order}
At the lowest threshold of each database, the baseline's mining call allocates 3.3 to
5.7 GiB, mostly for the new tid-list of each intersection. The optimized variant stores
each tidset as a bitset, a bit string in which bit $i$ stands for transaction $i$, made of
$W = \lceil N/64 \rceil$ 64-bit words. An intersection is a word-by-word AND, and the
support is the number of set bits, counted with \texttt{count\_ones} in the same loop.
The result goes into a bitset taken from a reusable pool. The bitset of an infrequent
intersection goes back to the pool, so a failed intersection allocates no memory. An
intersection reads $2W$ words whatever the support: $W = 50$ on Chess and $5\,316$ on
Accidents.

Items are visited in increasing order of support. The rare item that comes first is
intersected with the most items, but most of those intersections fail, and the frequent
item that comes last has few members after it, so child classes are smaller. The proof
in Section~\ref{sec:correctness} only needs a total order on the items, so it holds for
this order too.

\subsubsection{Perfect-extension pruning and single allocation of the result}
An item $b$ is a perfect extension~\cite{borgelt2012} of an itemset $X$ when
$t(X \cup \{b\}) = t(X)$, that is, when every transaction that contains $X$ also contains
$b$. Since $t(X \cup \{b\}) \subseteq t(X)$, the condition is equivalent to the two
supports being equal, so a single count comparison decides it. Every itemset $X \cup Y$
in the subtree of $X$ then has the same support as $X \cup Y \cup S$, where $S$ is any
subset of the set $E$ of perfect extensions. The optimized variant therefore removes $b$
from the child class and never intersects $b$'s tidset again. Each node of the search
tree stands for $2^{|E|}$ itemsets with the same support. On Mushroom 0.05 the search tree
has only 33\,819 nodes for 3\,755\,511 itemsets.

The search records only the nodes, in preorder. Each node holds its depth, the item it
adds, the perfect extensions it found, and its support. The function
\texttt{expand\_nodes} computes the final size of the result from the nodes, with an
overflow check, allocates the three result arrays exactly once, and then writes the
subsets $S$ in Gray-code order, so that two consecutive itemsets differ by one item. If
the result arrays grow one itemset at a time, each doubling of capacity copies the whole
array. On Mushroom 0.05, the mining call in an already warm process allocated 1\,068 MiB
to build a 306 MiB result and took 0.40 s. With a single allocation it allocates 310 MiB
and takes 0.17 s.

\subsubsection{Sparse data: pair counting, tid lists, and block bounds}
On sparse data, a bitset intersection still reads all $W$ words even when the tidsets are
short ($W = 1\,378$ on Retail and $1\,563$ on T10I4D100K), and most intersections fail.
On Retail 0.001, 99.9\% of the 2\,239\,786 intersections in the root class miss the
threshold. On T10I4D100K 0.00025, most intersections happen at depth 2 (2\,256\,822 of
them, 99.1\% failing) on tidsets of only about 100 tids. Three techniques address these
two patterns.

\emph{Pair counting at the root.} The function \texttt{count\_pairs} counts the support
of every pair of frequent items in a triangular array of $m_f(m_f-1)/2$ cells during one
pass over the horizontal database, as in Section 5.1 of the paper~\cite{zaki2000}, where
$m_f$ is the number of frequent items. The root class intersects only the pairs that
reach the threshold. At the root the prefix is empty, so the count of pair $(a, b)$ is
exactly $|t(\{a, b\})|$ and no frequent intersection is skipped. The counting runs only
when its estimated cost is below one tenth of the cost of the root intersections, that
is, when $10 \sum_i |T_i|(|T_i| - 1)/2 < W m_f(m_f-1)/2$. Retail and T10I4D100K meet this
condition. Chess, Mushroom, and Accidents, with long transactions and few frequent items,
intersect at the root as before.

\emph{Per-class tid lists.} A child class with $k$ members switches from bitsets to
increasing tid lists when
\begin{equation}
  kW + c\,(k - 1) \sum_{i=1}^{k} \sigma_i < W\,\frac{k(k-1)}{2},
  \label{eq:lists}
\end{equation}
where the right-hand side is the number of words the class's bitset intersections would
read, $kW$ is the cost of the conversion, $\sigma_i$ is the support of the $i$-th
member, and $c = 8$ is the cost of one merge step in words, chosen from a sweep of $c$
from 2 to 16 on the measurement machine. Rule~\eqref{eq:lists} depends on tidset lengths
and $N$, not on the name of the database. The merge of two lists has no data-dependent
branches, and it stops early once the common tids found so far plus the rest of the
shorter list cannot reach $\mathit{minsup}$.

\emph{Block bounds.} On T10I4D100K 0.00025, 98.7\% of the 2\,076\,972 merges still fail,
on two lists that hold 126 tids between them on average. Before it searches a class of
tid lists, the program splits the tid range into at most $K = 256$ blocks whose width is a
power of two, and counts the tids of each member in each block. For two members $X$ and
$Y$ of the class,
\begin{equation}
  \sigma(X \cup Y) \le \sum_{j=1}^{K} \min(c_{X,j}, c_{Y,j}),
  \label{eq:bound}
\end{equation}
where $c_{X,j}$ and $c_{Y,j}$ are the numbers of tids of $X$ and $Y$ in block $j$,
because a common tid lies in exactly one block and is counted once for each member. A
merge whose bound is below $\mathit{minsup}$ is skipped without reading any tid.
Bound~\eqref{eq:bound} never rejects a frequent intersection, so the output does not
change. We chose $K = 256$ from a sweep of 64 to 512 blocks.

\subsubsection{Start-up, reading, and writing}
The changes in this subsection are in the code shared by both variants. At high
thresholds, the runtime is not spent in the algorithm. One \texttt{demo.jl} process of
the baseline on Chess 0.9 takes 2.36 s, while Julia start-up takes only 0.17 s. The rest
is just-in-time (JIT) compilation. We run the public functions and the command-line
handling on a small database while Julia builds the package image, so their machine code
is stored with the package. The command-line handling lives in the package
(\texttt{src/cli.jl}) and not in the script file, because script code is not stored in
the package image and would be recompiled for 108.6 ms on every run. A test runs
\texttt{demo.jl} in a fresh process with \texttt{--trace-compile} and requires that no
method needs compiling.

At low thresholds on dense data, most of the time goes into writing the result. The
\texttt{Profile} sampler on Mushroom 0.05 showed that the comparison sort of the output
took 64\% of the time, formatting and writing the file 25\%, and mining only 11\%. The
optimized variant visits items by support, so it has to sort the result back into
lexicographic order. We replaced that sort with a radix sort on a 128-bit key that holds
the first items of the itemset and its index. The writer takes the digits of every number
below $2^{20}$ from a lookup table and writes them with one 8-byte store. The reader parses
each line into a reused vector and copies it once at its final length, so reading
Accidents (34 MB) in a warm process dropped from 0.57 s to 0.29 s and allocates 151 MiB
instead of 354 MiB.

\subsubsection{Julia features}
The type \texttt{ClassMember\{T\}} has a parameter \texttt{T}, which is \texttt{UInt64}
for bitsets or \texttt{Int} for tid lists. The search function \texttt{eclat\_search!}
is written once, and multiple dispatch (choosing a method from the types of all
arguments) picks the intersection, the conversion rule, and the intersection filter from
\texttt{T}. Julia compiles a specialized version for each tidset form, so the inner loop
has no run-time type checks. We use \texttt{Vector\{UInt64\}} instead of
\texttt{BitVector} because the intersection loop writes into reused memory and counts
bits in the same pass.

The key type \texttt{Tuple\{Vararg\{Int\}\}} of \texttt{SupportMap} has no fixed length,
so building millions of keys costs more than mining itself. In an early test on
Mushroom 0.05, inserting 3.76 million itemsets into the dictionary took 12.4 s, while
counting their supports took 0.63 s. Both variants therefore write their result into
three flat \texttt{Vector\{Int\}} arrays (the items, the start position of each itemset,
and the supports). The file writer reads these arrays directly, and \texttt{mine} builds
a \texttt{SupportMap} only at the interface boundary for the applications. Every tuning
constant is \texttt{const}, and every container on a hot path has a concrete element type.
\texttt{@inbounds} appears only in the tid-list merge and the number writer, with a
comment that proves the indices valid. On the bitset intersection loop,
\texttt{@inbounds @simd} changed mining time by $-8\%$ to $+11\%$ over five measurement
points, which is within noise, so we do not use it.

\subsection{Measuring the improvement}
\label{sec:measure}
Every number in this subsection is the median time of five cold \texttt{demo.jl}
processes: each run is a fresh process that includes Julia start-up, reading, mining, and
writing the file. The campaigns ran on a Linux machine (Intel Core i7-11800H, 7.4 GiB
RAM, clock fixed at 2.3 GHz) with single-threaded Julia 1.13.2. The grid has 33
thresholds on the five databases Chess, Mushroom, Retail, T10I4D100K, and Accidents,
plus some stress thresholds below the grid that only probe the limits. Besides SPMF's
Eclat, we compare with PrePost+~\cite{deng2015} from the
\href{https://github.com/NgTHung/Itemset/blob/main/tools/spmf/release.toml}{pinned SPMF release}
(Java Temurin 25.0.4.1, 4 GiB heap), because it is faster than SPMF's Eclat on dense
data. The ratio $R = T_\text{opt} / T_\text{PrePost+}$ divides the median time of the
optimized variant by the median time of PrePost+ at the same threshold.

Each technique was measured in its own session, in which the versions before and after
the change alternate on every repeat, so machine drift affects both equally. Only the
precompilation step was measured as two consecutive campaigns. Table~\ref{tab:steps} gives
the after-over-before time ratio of each step over the 33 thresholds, and the matching
peak-memory ratio. Peak memory is the peak resident set size (RSS), the largest amount of
physical memory the process holds, as measured by the operating system.

\begin{table*}[t]
\centering
\small
\caption{Effect of each optimization step over 33 thresholds: after-over-before ratios of
time and peak memory (below 1 is better), and the threshold that improved most}
\label{tab:steps}
\begin{tabular}{rllll}
\hline
Step & Technique & Time ratio & Example, time before $\to$ after (s) & Memory ratio \\
\hline
\optstep{step:precompile-package} & Precompile the package functions & 0.14 to 0.83 & Chess 0.9: 2.36 $\to$ 0.33 & 0.82 to 0.97 \\
\optstep{step:bitset} & Bitsets and increasing-support order & 0.10 to 0.97 & Mushroom 0.1: 6.28 $\to$ 0.88 & 0.68 to 0.99 \\
\optstep{step:perfect-extension} & Perfect-extension pruning & 0.93 to 1.04 & Mushroom 0.1: 0.89 $\to$ 0.83 & 0.95 to 1.00 \\
\optstep{step:radix-sort} & Radix sort for output order & 0.77 to 1.02 & Mushroom 0.1: 0.83 $\to$ 0.64 & 1.00 to 1.02 \\
\optstep{step:precompile-cli} & Precompile the command line & 0.66 to 0.94 & Chess 0.9: 0.33 $\to$ 0.22 & 0.91 to 1.03 \\
\optstep{step:pair-count} & Pair counting at the root & 0.24 to 1.01 & Retail 0.001: 1.72 $\to$ 0.42 & 1.00 to 1.03 \\
\optstep{step:tid-lists} & Per-class tid lists & 0.83 to 1.02 & T10I4D100K 0.00025: 1.46 $\to$ 1.21 & 1.00 to 1.01 \\
\optstep{step:byte-io} & Byte-level reading and writing, digit table & 0.76 to 0.99 & Accidents 0.65: 1.10 $\to$ 0.85 & 0.74 to 1.00 \\
\optstep{step:single-allocation} & Single result allocation, block bounds & 0.82 to 1.03 & T10I4D100K 0.00025: 1.18 $\to$ 0.97 & 0.83 to 1.02 \\
\hline
\end{tabular}
\end{table*}

Step~\ref{step:bitset} has the largest effect: time drops to 10\% to 97\% of the baseline
(median 37\%), and peak memory falls at every threshold. Step~\ref{step:pair-count} only
affects sparse data, where it cuts the time of Retail 0.001 by 76\%.
Step~\ref{step:perfect-extension} on its own shortens process time by at most 6.9\%,
because at that threshold sorting and writing hundreds of thousands of itemsets takes
most of the time. We keep it because it does not raise memory and because
step~\ref{step:single-allocation} builds on it. No step slows any threshold by more than
4.4\%.

Table~\ref{tab:final} shows the combined effect of steps~\ref{step:bitset}
to~\ref{step:single-allocation} at the highest and lowest threshold of each database. It
compares the baseline after step~\ref{step:precompile-package} with the optimized variant
after step~\ref{step:single-allocation}. Steps~\ref{step:precompile-cli}
and~\ref{step:byte-io} are in shared code, so the baseline of the submitted version also
benefits from them; the comparison of the two variants on the same code version is
step~\ref{step:bitset} of Table~\ref{tab:steps}. The two columns were measured in two
sessions three days apart on the same machine and configuration. PrePost+ ran in both
sessions, and its median differed between them by $-3.8\%$ to $+5.8\%$ (median
$+0.3\%$), so the two sessions are comparable within that margin.

\begin{table*}[t]
\centering
\small
\caption{Median cold process time (s) and maximum peak memory (MiB) over five runs. The
baseline (after step~\ref{step:precompile-package}) was measured on 2026-11-06, the
optimized variant (after step~\ref{step:single-allocation}) and PrePost+ on 2026-11-09,
and SPMF's Eclat on 2026-10-31, all on the same machine}
\label{tab:final}
\begin{tabular}{lrrrrrrrr}
\hline
 & & \multicolumn{5}{c}{Time (s)} & \multicolumn{2}{c}{Peak memory (MiB)} \\
Data / minsup & Itemsets & Baseline & Optimized & Speedup & SPMF Eclat & PrePost+ & Baseline & Optimized \\
\hline
Chess / 0.9 & 622 & 0.33 & 0.21 & 1.6 & 0.51 & 0.42 & 268 & $\le 254$ \\
Chess / 0.6 & 254\,944 & 3.24 & 0.30 & 10.8 & 15.98 & 0.59 & 308 & 262 \\
Mushroom / 0.5 & 153 & 0.33 & 0.22 & 1.5 & 0.51 & 0.44 & 256 & $\le 254$ \\
Mushroom / 0.1 & 574\,431 & 6.28 & 0.40 & 15.7 & 7.44 & 0.68 & 414 & 302 \\
Retail / 0.01 & 159 & 0.55 & 0.29 & 1.9 & 1.19 & 0.80 & 304 & $\le 254$ \\
Retail / 0.001 & 7\,589 & 12.80 & 0.38 & 33.5 & 1.53 & 3.52 & 331 & 281 \\
T10I4D100K / 0.01 & 385 & 2.93 & 0.32 & 9.1 & 0.78 & 1.42 & 304 & $\le 254$ \\
T10I4D100K / 0.00025 & 96\,959 & 10.92 & 0.97 & 11.3 & 8.02 & 4.34 & 363 & 279 \\
Accidents / 0.95 & 15 & 1.31 & 0.78 & 1.7 & 4.59 & 1.71 & 650 & 397 \\
Accidents / 0.65 & 1\,093 & 7.13 & 0.86 & 8.3 & 26.87 & 2.34 & 804 & 397 \\
\hline
\end{tabular}
\end{table*}

Across all 33 thresholds, steps~\ref{step:bitset} to~\ref{step:single-allocation} shorten
process time by 1.5 to 33.5 times (median 4.0). The geometric mean speedup per database is
3.1 on Accidents, 3.3 on Mushroom, 3.4 on Chess, 6.5 on Retail, and 13.5 on T10I4D100K.
The smallest gains are at the high thresholds of Chess and Mushroom, where a run takes
only about 0.22 s, of which 0.19 s is Julia start-up and package loading. Most of the gain
there comes from step~\ref{step:precompile-cli}. Against SPMF, the optimized variant is
2.3 to 53 times faster than Eclat, using the Eclat measurements from the 2026-10-31
session, and faster than PrePost+ at all 33 thresholds, with $R$ from 0.11 to 0.59. The
geometric mean of $R$ per database is 0.20 on Retail, 0.21 on T10I4D100K, 0.41 on
Accidents, 0.49 on Mushroom, and 0.51 on Chess. The interquartile range (IQR) of the
Baseline and Optimized columns stays within 3.3\% of the median at every threshold.

After step~\ref{step:single-allocation}, peak memory is lower than the baseline's at
every threshold. The largest drop is on Accidents (804 to 397 MiB at 0.65), because the
reader no longer over-allocates and the tidsets take only a few MiB during the search. On
the measurement machine, all three programs (Julia, SPMF Eclat, and PrePost+) report
their lowest values between 253.0 and 253.6 MiB, so the measurement has a common floor,
and the cells marked $\le 254$ in Table~\ref{tab:final} are only upper bounds.

We tried, measured, and dropped three techniques. Diffsets~\cite{zaki2003} store the tids
that a member lacks relative to its prefix. In a sweep of the switching ratio, no value
was faster than bitsets: Accidents 0.2 slowed by up to 10\% and T10I4D100K 0.00025 by up
to 71\%. The reason is that perfect-extension pruning already removes every member with
an empty diffset, which is exactly where diffsets help most. Even with the best choice for
each class, the number of words read dropped by at most 0.3\%. Printing itemsets directly
from the sort key was 5\% to 7\% slower. The \texttt{@inbounds @simd} annotation on the
bitset intersection gave no improvement beyond noise.

The optimized variant has three measured limits. First, at the stress threshold
Mushroom 0.05, the program has to write 3\,755\,511 lines and remains slower than
PrePost+ (1.45 s against 1.38 s). Second, each process spends about 0.19 s on Julia
start-up and package loading, which could only be removed by compiling the whole program
ahead of time into an executable. Third, pair counting raises peak memory by 20\% at the
stress threshold Retail 0.00025, because the triangular array for 5\,918 frequent items
takes 67 MiB.
